\begin{definition}[Measure root consumer package]
\label{def:measure-root-consumer-package}
The $\MeasureUp$ root consumer package is the root NameCert public surface of
\autoref{def:measure-root-namecert-public-surface} read together with the
downstream package of
\autoref{thm:measure-root-downstream-unblock-package}. Its consumers receive
only covered measurable-event rows, $\RealUp$ endpoint classifier rows,
countable-sum endpoint rows, head-tail readback, empty-zero, finite
disjoint-union readback, relative-difference additivity, monotone inclusion,
nonnegative endpoint, and the explicit no-$\StdBridge$ row exposed by that
public $\MeasureUp$ package.
\end{definition}

\begin{theorem}[Measure integral consumer rows]
\label{thm:measure-integral-consumer-rows}
The root consumer package of
\autoref{def:measure-root-consumer-package} supplies $\IntegralUp$ with
exactly the covered measurable-event rows, $\RealUp$ endpoint classifier rows,
countable-sum endpoint rows, head-tail readback rows, and finite
disjoint-union readback rows exposed by $\MeasureUp$. No host set equality,
quotient event equality, probability normalization row, functional-analysis
order row, or standard bridge row is added to the integral consumer surface.
\end{theorem}
\begin{proof}
Project the root NameCert public surface from
\autoref{def:measure-root-consumer-package}. Its coverage predicate is
\autoref{def:measure-event-row-coverage}, so every event row visible to the
consumer is one of the displayed measurable-event rows of the $\MeasureUp$
surface.

Apply \autoref{thm:measure-root-downstream-unblock-package}. This exposes the
covered event rows, $\RealUp$ endpoint rows, countable-sum rows, head-tail
readback, and finite disjoint-union readback as the public downstream package.

Now apply \autoref{thm:measure-integral-premeasure-consumption}. Its
$\IntegralUp$ projection consumes exactly those rows from the same public
package. Repacking that projection with the root consumer package gives the
claimed integral consumer surface, and the exactness predicate rules out every
row not named in the statement.
\end{proof}

\begin{theorem}[Measure probability consumer rows]
\label{thm:measure-probability-consumer-rows}
The root consumer package of
\autoref{def:measure-root-consumer-package} supplies $\ProbSpaceUp$ with the
total-event normalization row together with the empty-zero and finite
disjoint-union rows exposed by $\MeasureUp$. No integral countable-sum
endpoint row, functional-analysis relative-difference package, quotient event
equality, or standard bridge row is added to the probability consumer surface.
\end{theorem}
\begin{proof}
Project the empty-zero and finite disjoint-union rows from
\autoref{thm:measure-root-downstream-unblock-package}. These are already
public rows of the root consumer package, because the package includes the
root NameCert surface and its downstream repacking.

Apply \autoref{thm:measure-probability-normalization-row}. The theorem reads
the total-event endpoint and the unit-mass comparison through the same
$\MeasureUp$ public surface, and it consumes the empty-zero and finite
disjoint-union rows already projected above.

Repack the total-event normalization row with those two additive rows. The
event-row exactness clause of
\autoref{def:measure-root-consumer-package} leaves no additional event
identity, countable-sum endpoint, functional-analysis order, or
$\StdBridge$ row in the $\ProbSpaceUp$ consumer surface.
\end{proof}

\begin{theorem}[Measure functional-analysis consumer rows]
\label{thm:measure-functionalanalysis-consumer-rows}
The root consumer package of
\autoref{def:measure-root-consumer-package} supplies $\FunctionalAnalysisUp$
only with the nonnegative endpoint, monotone inclusion, relative-difference,
and countable-sum classifier rows exposed by $\MeasureUp$. It does not add
integral premeasure rows, probability normalization rows, quotient event
equality, host set equality, or a standard bridge row.
\end{theorem}
\begin{proof}
Start from the $\MeasureUp$ analysis substrate rows already present in the
root consumer package. The nonnegative endpoint, monotone inclusion,
relative-difference, and countable-sum classifier rows are the analysis rows
visible through the public root surface.

Apply \autoref{thm:measure-functionalanalysis-substrate-obligations}. Its
$\FunctionalAnalysisUp$ projection consumes exactly those four row families:
nonnegative endpoints, monotone inclusion, relative-difference additivity,
and countable-sum classifier stability.

Finally repack the projection through
\autoref{thm:measure-root-downstream-unblock-package}. The downstream
exactness predicate keeps the consumer boundary at the named row families, so
the functional-analysis surface gains neither integral rows, probability
normalization rows, quotient event equality, host set equality, nor a
$\StdBridge$ row.
\end{proof}
